\section{Related Work}
\label{sec:related}

\subsection{Decision Tree Induction: CART and C4.5}

CART~\cite{breiman1984cart} builds a binary tree by recursive partitioning. At each node, it chooses the split that most reduces Gini impurity. After growth, it simplifies the tree with cost-complexity pruning. C4.5~\cite{quinlan1993c45} chooses splits by gain ratio, defined as information gain divided by split entropy. It creates one branch per value for nominal attributes, splits numeric attributes at a single threshold, and applies error-based pruning. For a numeric attribute, a full search sorts the values in the node and evaluates every threshold between consecutive distinct values. For entropy-based criteria, Fayyad and Irani~\cite{fayyad1992continuous} showed that an optimal threshold lies on a boundary between classes in the sorted order, allowing the remaining thresholds to be skipped without changing the result. EC4.5~\cite{ruggieri2002yadt} produces the same tree as C4.5 with faster sorting and threshold search, and its successor YaDT~\cite{ruggieri2004yadt} further restricts the search to these class boundaries. In both algorithms, sorting and scanning the samples at each node dominate the training time.

\subsection{Parallel Induction on CPUs}

SLIQ~\cite{mehta1996sliq} sorts each numeric attribute once before induction and grows the tree breadth-first. It keeps the class labels in a list that must fit in memory. SPRINT~\cite{shafer1996sprint} removes this memory-resident class list by storing class labels inside each presorted attribute list. When a node is split, each list is partitioned between the children rather than re-sorted. SPRINT also runs on clusters by distributing the lists across machines. Both systems still evaluate every distinct threshold, scaling by avoiding repeated sorting rather than by approximating the split search.

Tree induction can be parallelized through task parallelism, which grows different subtrees concurrently, or data parallelism, which divides work within a single node across attributes or samples. Because subtrees vary widely in size, task parallelism alone suffers from load imbalance, whereas data parallelism runs out of work once nodes become small. Many systems therefore combine both, using data parallelism near the root and task parallelism deeper in the tree. Aldinucci et al.~\cite{aldinucci2014fastflow} parallelized YaDT on multicore processors with the FastFlow framework: their NP strategy processes nodes as tasks, while their NAP strategy also splits a node's attributes across threads when a cost model predicts a benefit. They call this combination ``hybrid'', though it targets only multicore CPUs. \plan{Verify ``hybrid'' in the CCPE 2014 PDF (the 2010 version does not use the word); drop the quote if absent.} Anghel et al.~\cite{anghel2019bfdn} train random forests on multicore CPUs with an exact tree builder that grows the upper levels breadth-first and small nodes depth-first. \plan{Verify in the arXiv PDF how they switch from breadth-first to depth-first (node size, cache fit).} Choosing the parallelization strategy by node size mirrors our choice of execution device.

Yggdrasil~\cite{abuzaid2016yggdrasil} partitions data by attribute across a Spark cluster so each worker returns only its best split. This avoids the binning that sample-partitioned distributed systems need to limit communication, but Yggdrasil targets clusters rather than a single machine.

\subsection{GPU-Based Tree Induction}

\paragraph{Single Trees.} CUDT~\cite{lo2014cudt} ports SPRINT's attribute lists to CUDA. The CPU controls the induction and stores the tree. The GPU evaluates the splits with prefix sums and partitions the lists, one node at a time. Its speedup is largest at the root and shrinks with depth as nodes become smaller and more numerous, each requiring host--device transfers. Nasridinov et al.~\cite{nasridinov2014gpu} divide the work similarly for ID3, growing one node at a time: the GPU sorts the node's samples and evaluates candidate splits, while the CPU updates the tree. They describe this as ``hybrid CPU--GPU'' and report both runtime and energy on a low-end GPU. \plan{Paywalled, verified from the abstract only. Check in the PDF: GPU sorts and evaluates splits, CPU updates the tree; the quote ``hybrid CPU--GPU''; the GPU model (GT220?).} Strnad and Nerat~\cite{strnad2016gpu} instead keep the entire induction of a CART-like binary tree on the GPU, combining node-, attribute-, and split-level parallelism. Their scheduler groups nodes of similar size into batches to reduce load imbalance, but still runs every batch on the GPU.

\paragraph{Trees in Ensembles.} Several GPU systems grow trees as part of ensembles, and their per-tree techniques apply to single trees. CudaTree~\cite{liao2013cudatree} trains random forests, growing large nodes depth-first and switching to breadth-first growth below a size threshold, all on the GPU. In its hybrid mode, the GPU and a multicore CPU library (scikit-learn or wiseRF) train one forest together, but each device trains its own trees. \plan{CudaTree paper PDF not found (link dead); details checked against the CudaTree GitHub README only.} For gradient boosting, Mitchell and Frank~\cite{mitchell2017xgboostgpu} grow each XGBoost~\cite{chen2016xgboost} tree entirely on the GPU, level by level, using prefix sums over sorted attribute values. ThunderGBM~\cite{wen2019thundergbm} likewise identifies the shift from a few large nodes to many small ones as a main challenge on GPUs. On the GPU, current versions of XGBoost and cuML~\cite{raschka2020cuml}, which we compare against, use histogram-based split finding instead.

\subsection{Heterogeneous CPU--GPU Induction}

In the systems above, either the CPU only controls induction while the GPU computes splits (CUDT, Nasridinov et al.), size-aware scheduling stays on the GPU (CudaTree, Strnad and Nerat), or the two devices train different trees of a forest (CudaTree's hybrid mode). A patent by Ioannou et al.~\cite{ioannou2025patent} divides the nodes of a tree between devices: the GPU builds large sibling nodes level by level, while the CPU concurrently grows small ones depth-first. It is aimed mainly at tree ensembles and prefers histogram-based split finding. Our hybrid scheme uses the same division for exact CART and C4.5, with a pool of CPU threads growing the subtrees below a size threshold.
